\documentclass[UTF8,fontset=windows,zihao=-4]{ctexart}

\usepackage[a4paper,left=22mm,right=22mm,top=24mm,bottom=22mm,headheight=18pt]{geometry}
\usepackage{graphicx}
\usepackage{array}
\usepackage{booktabs}
\usepackage{longtable}
\usepackage{tabularx}
\usepackage{enumitem}
\usepackage{fancyhdr}
\usepackage{hyperref}
\usepackage{caption}
\usepackage{float}
\usepackage{setspace}
\usepackage{xcolor}

\graphicspath{{figures/}}
\hypersetup{hidelinks}
\captionsetup{font=small,labelfont=bf}
\setstretch{1.35}
\setlength{\parindent}{2em}
\setlength{\parskip}{0.2em}
\setlist{nosep,leftmargin=2.2em}
\renewcommand{\arraystretch}{1.25}

\newcommand{\softwareName}{@@software_name@@}
\newcommand{\softwareVersion}{@@software_version@@}
\newcommand{\copyrightOwner}{@@copyright_owner@@}

\newcommand{\menu}[1]{\textbf{#1}}
\newcommand{\api}[1]{\nolinkurl{#1}}
\newcommand{\redblank}[1]{{\color{red}\underline{\hspace{#1}}}}

\pagestyle{fancy}
\fancyhf{}
\fancyhead[L]{\softwareName\ \softwareVersion}
\fancyhead[R]{第 \thepage 页}
\fancyfoot[L]{著作权人：\copyrightOwner}
\fancyfoot[R]{使用说明书}
\renewcommand{\headrulewidth}{0.4pt}
\renewcommand{\footrulewidth}{0.4pt}

\begin{document}

\begin{titlepage}
  \centering
  \vspace*{28mm}
  {\zihao{1}\bfseries \softwareName\par}
  \vspace{8mm}
  {\zihao{2}\bfseries \softwareVersion\ 使用说明书\par}
  \vspace{18mm}
  {\zihao{4} 计算机软件著作权登记文档鉴别材料\par}
  \vfill
  \begin{tabular}{>{\bfseries}p{36mm}p{86mm}}
    软件全称 & @@software_name@@ \\
    软件版本 & \softwareVersion \\
    开发完成日期 & @@development_date@@ \\
    文档类型 & 使用说明书 \\
    适用对象 & @@audience@@ \\
  \end{tabular}
  \vfill
  {\zihao{-4} 本说明书适用于 \softwareName\ \softwareVersion，说明软件功能组成、运行环境、操作流程、主要界面和数据接口。\par}
\end{titlepage}

\begingroup
\small
\setstretch{1}
\setlength{\parskip}{0pt}
\tableofcontents
\endgroup
\clearpage


% Replace every pending paragraph with facts verified against the NEW project's real code and screenshots.
\section{软件概述}
待编写：研发目的、实际功能、使用对象和适用范围。
\section{运行环境与安装启动}
待编写：操作系统、运行时、依赖、启动命令及停止、备份方法。
\section{主要界面与操作流程}
待编写：实机截图、逐步操作、输入、反馈及结果。截图放入figures目录。
% Example syntax (uncomment only after capturing the real screen):
% \begin{figure}[!htbp]\centering
% \includegraphics[width=0.95\linewidth]{figures/01-home.png}
% \caption{实际软件主界面}\end{figure}
\section{结果查询与数据导出}
待编写：仅描述真实实现的报告、记录、导出或查询功能。
\section{异常处理及常见问题}
待编写：真实前置条件、失败提示、纠错步骤和边界。
\section{数据保存与软件结构}
待编写：存储、备份、访问范围以及与当前版本一致的软件结构。
\end{document}
